\begin{python0}
from solutions import *; clear()
\end{python0}
\sectionthree{Summary}

A function looks like this:

\texttt{ [return type] [func name]}( \texttt{[param type] [var name]} )

\{

\texttt{[body of function]}

\}

There can be more than one parameter type and name. A parameter is just a variable that receives a value passed into the function. The return type can be any valid type including \texttt{void} if the function does not return a value.

The function body can be any valid C++ statement. To return a value the
function executes

\texttt{return [expression];}

in its body. If the function does not return a value, executing
\begin{center}
\texttt{return;}
\end{center}
will return to the point of function call.

If a function has return type void, return need not be present. In that case when the execution reached the end of the function block, a return is executed.

Functions cannot be nested.

The largest scope in a program file is the global scope. You can declare variables and constants in the global scope. These are called global variables and global constants respectively. However you should avoid (if possible) the declaration of global variables.

On returning from a function, since the program leaves the scope of the function, all variables created within the function's scope are destroyed.

A parameter is pass-by-value if the parameter has its own value. A parameter is pass-by-reference if the parameter does not have its own value but references the value of a variable that is passed to that parameter. Changing the value of the parameter that is pass-by-reference will change the value of the variable passed to the parameter. By default, parameters of basic types (int, double, char, bool) are pass-by-value. The following parameter \texttt{x} is pass-by-value:
\begin{center}
\texttt{void f(..., int x, ...) { ... }}
\end{center}
while the following parameter \texttt{y} is pass-by-reference:
\begin{center}
\texttt{void g(..., int \& x, ...) { ... }}
\end{center}
A function prototype is just the header of a function as a statement. A function prototype is usually placed near the top of source files in order to indicate to the compiler the existence that function. If a function prototype is seen by the compiler, then calling that function will be compiled correctly by the compiler even if the function definition is below the point where the function call is made.

